\documentclass{article}
\fontsize{10.5}{12}\selectfont
% Language setting
% Replace `english' with e.g. `spanish' to change the document language
\usepackage[english]{babel}

% Set page size and margins
% Replace `letterpaper' with `a4paper' for UK/EU standard size
\usepackage[a4paper,top=3.5cm,bottom=3.4cm,left=3cm,right=3cm]{geometry}

\usepackage{bold-extra}
% \usepackage{CJKutf8}
% \usepackage[utf8x]{inputenc}
% \newcommand\customfont[1]{{\usefont{T1}{msgothic}{m}{n} #1 }}

% \usepackage{xeCJK}
% \setCJKmainfont{IPAMincho} % for \rmfamily
% \setCJKsansfont{IPAGothic} % for \sffamily
% \usepackage{xltxtra,fontspec,xunicode}
% \setmainfont{MSGOTHIC}
\usepackage{fontspec}
\usepackage{zxjatype}
\usepackage[ipa]{zxjafont}
% Section Heading Spacing
\usepackage{titlesec}
\titlespacing*{\section}{0pt}{0.1em}{-0.1em}
\titlespacing*{\subsection}{0pt}{0.1em}{-0.1em} 
\titleformat*{\section}{\bfseries}
\titleformat*{\subsection}{\bfseries}
% \setlength\intextsep{6pt}
% \setlength\itemsep{-15.5em}% item spacing
%list
\usepackage{enumitem}
\setlist[itemize]{leftmargin=*}


\pagenumbering{gobble} %eliminate page number
\newcommand*{\tabindent}{ \hspace{3mm}} %table indent
\usepackage{anyfontsize}
% \renewcommand{\arraystretch}{1.45}
\usepackage{array}
\newenvironment{mytable}[1][]{\begin{table}[#1]\vspace{1.0cm}}{\vspace{1.0cm}\end{table}}
 \usepackage{parskip}

% Useful packages
\usepackage{amsmath}
\usepackage{graphicx}
\usepackage[colorlinks=true, allcolors=blue]{hyperref}
% \usepackage[caption=false, font=normalsize, labelfont=sf, textfont=sf]{subfig}
\usepackage[caption=false, font=footnotesize]{subfig}
\usepackage{float}
\newcommand\nocell{\multicolumn{1}{c}{}}
\usepackage{multicol}
\setlength{\columnsep}{0.75cm}
\usepackage{cite}
\def\citepunct{,\,}

\usepackage{lipsum}

\title{}
\author{}
% \setCJKmainfont{IPAMincho}
\setlength\extrarowheight{5pt}

\begin{document}

\def \mysection#1{
\section{\large{#1}}
}

\fontsize{12pt}{19pt}
% \maketitle
\begin{center}
\setmainfont{Arial}[BoldFont=IPAGothic]
\textbf{\\ 社会文化環境に関する研究}\\
\fontsize{10.8}{18}\selectfont
Study on Socio-Cultural Environments %\newline
\end{center}
\begin{table}[h]
  \begin{center}
  \vspace{0.3cm}
  \begin{tabular}{ m{1.83cm} m{5.35cm} }
        \nocell{学籍番号} & \tabindent \fontsize{11}{10}\selectfont \textsc{47201111} \\%[5pt]
        \nocell{氏　　名} & \tabindent 環境 \enspace 馨\fontsize{11}{10}\selectfont（Kankyo, Kaoru）\\%[5pt]
        \nocell{指導教員} & \tabindent 環境\enspace 太郎\enspace 教授\\%[5pt]
  %\multicolumn{4}{l}{$^{\mathrm{a}}$Sample of a Table footnote.}
  \end{tabular}
  \end{center}
%   \vspace{-0.2cm}
\end{table}

\begin{multicols}{2}
\setmainfont{Times New Roman}

\fontsize{10.5}{19}\selectfont
% \begin{abstract}
% Your abstract.
% \end{abstract}
%%%%%%%%%%%%%%%%%%%%%%%%%%%%%%% Content Begins here %%%%%%%%%%%%%%%%%%%%%%%%%%%%%%%%%%%%%%%%%%%%%%%%%%%%%%%%%%%%%%%%%%%%%%%%%%%%%%%%%
%%%%%%%%%%%%%%%%%%%%%%%%%%Please use \begin{figure}[H] or \begin{table}[H] for figure and table, figure* table* is fine%%%%%%%%%%%%%%%%%%
ここから本文を書き始める．

上35mm，下30mm，左右それぞれ30mmの余白を確保する．

1ページ目は1行空けて論文題目を記入し，その下に英訳を記入する．

次に1行空けて，学籍番号，氏名，および指導教官を記入する．
さらに1行空けて本文を書きはじめる．

ページはこちらで打つので記入しないこと．
ただし，各ページ裏面に鉛筆書きで氏名およびページ番号を記入すること．

テキストは2段組とする．図表等の配置で2段組にならない部分は適宜判断してよい．

フォント

論文題目12ポイントのゴチック体
英文の場合、Arial or Bold Face
本文は10.5ポイントの明朝体
英文の場合、Times New Roman

図表を含めて片面4ページとする．

中間発表(9月修了者は4月，3月修了者は10月)の時には会場配布用に20部印刷して教員には個別に配布する．

2月の最終発表会(3月修了者)の時にはこの原稿を印刷した「梗概集」が配られるので会場配布用に印刷する必要はない．

8月の最終発表会(9月修了者)に際しては、中間発表時と同様に20部印刷して教員には個別に配布する．

このフォーマットはサイボウズにある図書館公開用の要旨とは若干フォーマットが異なっているが，図書館公開用CD-Rに収録するファイルもこちらの様式で作成する\cite{einstein}．

\lipsum[1-10]
% \input{01_Introduction}
% % \input{02_RelatedWork}
% \input{03_Motivation}
% \input{05_Approach}
% \input{07_Evaluation}
% % \input{08_Discussion and Future Works}
% \input{09_Conclusion}



\bibliographystyle{ieeetr}
\bibliography{sample}

\end{multicols}

\end{document}